% ---------------------------------------------------------------------
% 2.5 Interfaces cerebro-computadora
% ---------------------------------------------------------------------
\section[INTERFACES CEREBRO-COMPUTADORA]{Interfaces cerebro-computadora}
\label{sec:bci}

\subsection{Definición y componentes}

Las interfaces cerebro-computadora (BCI) permiten entablar comunicación directa entre la
actividad cerebral de un individuo y un dispositivo externo, sin depender de ninguna otra señal o
movimiento muscular. De acuerdo con Wolpaw et al., una BCI es un sistema de comunicación en el que
los mensajes o comandos que el usuario envía al mundo exterior no pasan por las vías de salida
normales del cerebro, es decir, por los nervios periféricos y los músculos \cite{wolpaw2002}, y
como se menciona en el trabajo de Felton et al., estos sistemas traducen las señales cerebrales
registradas en acciones en tiempo real, lo cual abre la posibilidad de restaurar funciones
comunicativas o motoras en personas con discapacidad severa \cite{felton2007}.

Independientemente de la señal que utilice, una BCI se compone de una serie de etapas que se
ejecutan de forma cíclica, como se ilustra en la Figura~\ref{fig:bci}. Primero se adquiere la señal
cerebral y se digitaliza; después se preprocesa para atenuar el ruido y los artefactos; a
continuación se extraen las características que mejor representan la intención del usuario y un
algoritmo de clasificación las traduce en un comando para el dispositivo, cuyo resultado se
presenta al usuario como retroalimentación para que pueda ajustar su estrategia
\cite{wolpaw2012,nicolasalonso2012}.

\begin{figure}[htbp]
\centering
\begin{tikzpicture}[
  node distance=0.38cm,
  etapa/.style={draw, rounded corners=2pt, fill=gray!8, align=center, minimum height=1.1cm, text width=2.15cm, font=\footnotesize},
  flecha/.style={-{Stealth[length=2mm]}, thick}]
  \node[etapa, fill=red!8] (usr) {Usuario\\(actividad cerebral)};
  \node[etapa, right=of usr] (adq) {Adquisición\\(EEG)};
  \node[etapa, right=of adq] (pre) {Preprocesa-\\miento};
  \node[etapa, right=of pre] (car) {Extracción de\\características};
  \node[etapa, right=of car] (cla) {Clasificación};
  \node[etapa, below=1.1cm of cla, fill=blue!6] (app) {Aplicación\\(deletreador)};
  \draw[flecha] (usr) -- (adq);
  \draw[flecha] (adq) -- (pre);
  \draw[flecha] (pre) -- (car);
  \draw[flecha] (car) -- (cla);
  \draw[flecha] (cla) -- node[right, font=\footnotesize] {comando} (app);
  \draw[flecha, dashed] (app.west) -| node[pos=0.25, below, font=\footnotesize] {retroalimentación visual} (usr.south);
  \begin{scope}[on background layer]
    \node[draw=gray, dashed, rounded corners, fit=(pre)(car)(cla), inner sep=6pt,
          label={[font=\footnotesize]above:Procesamiento de la señal}] {};
  \end{scope}
\end{tikzpicture}
\caption{Diagrama general de una interfaz cerebro-computadora. Elaboración propia con base en
\cite{wolpaw2002,wolpaw2012}.}
\label{fig:bci}
\end{figure}

\subsection{Antecedentes históricos}

El término interfaz cerebro-computadora fue introducido por Jacques Vidal en 1973, quien planteó
por primera vez la posibilidad de utilizar las señales del EEG como un canal de comunicación con
una computadora \cite{vidal1973}. Quince años después, Farwell y Donchin presentaron el primer
deletreador basado en la componente P300 \cite{farwell1988}, y en 1999 Birbaumer et al.
demostraron que pacientes con ELA podían escribir mensajes controlando potenciales corticales
lentos \cite{birbaumer1999}. A partir de la primera reunión internacional sobre BCI, celebrada en
1999 \cite{wolpaw2000}, se organizaron las competencias internacionales de BCI, cuyas bases de
datos públicas permitieron comparar algoritmos bajo las mismas condiciones y siguen siendo una
referencia obligada para la detección de la P300 \cite{blankertz2004,blankertz2006}. Durante la
última década, el uso de redes neuronales profundas ha desplazado a los métodos basados en
características diseñadas manualmente \cite{craik2019,roy2019}, y recientemente se han incorporado
modelos de lenguaje de gran escala para acelerar la escritura \cite{hong2026}. La
Figura~\ref{fig:historia} resume estos hitos.

\begin{figure}[htbp]
\centering
\begin{tikzpicture}[
  hito/.style={align=center, font=\footnotesize, text width=2.2cm},
  punto/.style={circle, fill=red!60!black, inner sep=2pt}]
  \draw[thick, -{Stealth[length=2mm]}] (0,0) -- (15.2,0);
  \foreach \x/\anio/\texto/\pos in {
      0.4/1929/{Berger: primer EEG en humanos \cite{berger1929}}/above,
      2.1/1965/{Sutton et al.: descubrimiento de la P300 \cite{sutton1965}}/below,
      3.8/1973/{Vidal: concepto de BCI \cite{vidal1973}}/above,
      5.5/1988/{Farwell y Donchin: P300 Speller \cite{farwell1988}}/below,
      7.2/1999/{Birbaumer et al.: deletreo con ELA \cite{birbaumer1999}}/above,
      8.9/2004--06/{Competencias BCI II y III \cite{blankertz2004,blankertz2006}}/below,
      10.6/2011/{CNN para detección de P300 \cite{cecotti2011}}/above,
      12.3/2018/{EEGNet \cite{lawhern2018}}/below,
      14.0/2026/{P300 Speller con LLM \cite{hong2026}}/above}{
    \node[punto] at (\x,0) {};
    \node[font=\footnotesize\bfseries, \pos=2pt] at (\x,0) {\anio};
    \ifthenelse{\equal{\pos}{above}}
      {\node[hito, above=0.45cm] at (\x,0) {\texto};}
      {\node[hito, below=0.45cm] at (\x,0) {\texto};}
  }
\end{tikzpicture}
\caption{Línea de tiempo de algunos hitos en el desarrollo de las BCI basadas en EEG. Elaboración
propia.}
\label{fig:historia}
\end{figure}

\subsection{Clasificación de las BCI}

Las BCI pueden clasificarse bajo distintos criterios. Según la forma en que se registra la señal,
se dividen en invasivas, que requieren colocar electrodos dentro o sobre el cerebro, y no
invasivas, que registran la actividad desde el exterior \cite{nicolasalonso2012}. Según el origen
de la actividad que aprovechan, se distinguen las exógenas, que dependen de una estimulación
externa como la iluminación de una matriz de caracteres, y las endógenas, en las que el usuario
modula voluntariamente su actividad cerebral sin estímulo alguno, por ejemplo al imaginar el
movimiento de una mano \cite{nicolasalonso2012}. Así mismo, una BCI es síncrona cuando el sistema
determina los intervalos en los que el usuario puede emitir un comando, y asíncrona cuando el
usuario puede hacerlo en cualquier momento, lo que exige detectar también los periodos en los que
no desea comunicarse \cite{wolpaw2012}. Bajo estos criterios, el P300 Speller es una BCI no
invasiva, exógena y síncrona.

\subsection{Paradigmas de control}

Las BCI basadas en EEG utilizan principalmente cuatro tipos de señales, cuyas características se
comparan en la Tabla~\ref{tab:paradigmas}. Las basadas en ritmos sensoriomotores aprovechan la
atenuación de los ritmos mu y beta durante la imaginación motora; las de potenciales evocados
visuales de estado estable (SSVEP) utilizan la respuesta de la corteza visual a estímulos que
parpadean a una frecuencia fija; las de potenciales corticales lentos requieren que el usuario
aprenda a desplazar voluntariamente la línea base de su EEG; y las basadas en la P300 aprovechan
la respuesta del paradigma del evento raro descrita en la sección anterior
\cite{nicolasalonso2012,wolpaw2002}. Una de las principales ventajas de este último enfoque es que
la P300 es una respuesta natural ante los estímulos relevantes, por lo que prácticamente no
requiere entrenamiento previo del usuario \cite{farwell1988,contreras2021}.

\begin{table}[htbp]
\centering
\caption{Comparación de los principales paradigmas de control de BCI basadas en EEG}
\label{tab:paradigmas}
\small
\begin{tabularx}{\textwidth}{L{2.2cm}L{2.6cm}L{2.5cm}YY}
\toprule
\textbf{Paradigma} & \textbf{Señal} & \textbf{Entrenamiento del usuario} & \textbf{Ventajas} & \textbf{Desventajas} \\
\midrule
Imaginación motora (SMR) & Ritmos mu y beta sobre corteza motora & Semanas & Endógena; no requiere estímulos externos. & Pocas clases; alta variabilidad entre usuarios. \\
SSVEP & Respuesta a estímulos que parpadean a frecuencia fija & Mínimo & Alta tasa de transferencia; señal robusta. & Fatiga visual; requiere control de la mirada. \\
Potenciales corticales lentos (SCP) & Desplazamientos lentos de la línea base & Semanas a meses & Funcionó con pacientes con ELA \cite{birbaumer1999}. & Muy lenta; entrenamiento prolongado. \\
P300 & Componente P300 del paradigma \textit{oddball} & Mínimo & Muchas opciones (p.\,ej. 36 caracteres); respuesta natural. & Requiere repeticiones; sensible a la fatiga. \\
\bottomrule
\end{tabularx}
\par\smallskip{\footnotesize Elaboración propia con información de \cite{nicolasalonso2012,wolpaw2002,lotte2018}.}
\end{table}

\subsection{Calibración y variabilidad entre sujetos}

La forma de onda de la P300 varía considerablemente de una persona a otra e incluso entre sesiones
de una misma persona, por lo que la mayoría de los sistemas requiere una etapa de calibración en la
que el usuario deletrea un texto conocido para entrenar un clasificador específico para él
\cite{wolpaw2012}. Esta dependencia ha motivado el estudio de esquemas de clasificación entre
sujetos, en los que el modelo se entrena con datos de otras personas; en este escenario,
Alvarado-González et al. reportaron que el desempeño de las redes convolucionales disminuye
respecto al entrenamiento intrasujeto, aunque la diferencia no es particularmente grande
\cite{alvarado2021}. Por otra parte, no todas las personas logran controlar una BCI con la misma
facilidad: en un estudio con 100 participantes, Guger et al. encontraron que el 72.8\% de los 81
sujetos que utilizaron el paradigma de filas y columnas deletreó con una exactitud del 100\%
empleando apenas alrededor de cinco minutos de datos de calibración, mientras que con el
paradigma de carácter único esta proporción bajó a 55.3\% \cite{guger2009}.
